\begin{abstract}
Language-model systems can answer from parametric weights, retrieve external evidence, write reusable external memory, apply temporary task-time updates, or consolidate stable updates more durably. Existing work usually studies these operations in isolation or focuses only on retrieval selection. We present a unified multi-timescale routing framework that chooses among these operations under a single reward that trades off answer quality, latency, write cost, and forgetting risk while incorporating delayed utility. Our empirical story has two parts. On PopQA, the router matches retrieval-level quality overall after guardrails, preserves recurring-case quality, and materially reduces repeated retrieval on recurring queries; these gains depend on memory reading, memory writing, and future-value-aware routing. On a deterministic bundled FreshQA-style update benchmark, the router improves updated-knowledge accuracy and stale-answer rate over retrieval-only and memory-only baselines while using meaningful temporary adaptation, rare guarded durable consolidation, and explicit rollback. Together these results support the claim that memory operations should be routed jointly rather than optimized as isolated components.
\end{abstract}
